\chapter{The Shape of the Start}

\nopagebreak
\epigraph{"Moore's law and software bloat"}{Title of the first post, January~9,~2008 (182~words)}
\nopagebreak

On January 9, 2008, the first post appeared. It observed that software expands to consume whatever hardware speed Moore's law provides. One hundred eighty-two words. No code. No equations. Just a small observation.

A second post the same day: "Integration surprises." One hundred thirty-four words about the practical value of techniques that seemed impractical in college.

Cook posted twice on his first day. He posted again on January 10. And 11. And 12. By the end of January, there were 37 posts. By the end of 2008, there were 410, capturing the energy of beginning.

\bigskip

\section*{The Launch Year}

2008 stands alone in the archive. Four hundred and ten posts in a single year, more than one per day on average. The posts are short: 206 words on average, some under 100. They read like notes jotted between other work.

A quick observation about software bloat. A pointer to an interesting integral. A link to someone else's insight. A tip about PowerShell. A note on clinical trial methodology.

The topics scatter across a consultant's working life:

\begin{longtable}{l p{6.5cm}}
\toprule
\textbf{Category} & \textbf{Posts in 2008} \tabularnewline
\midrule
Software development & 86 \tabularnewline
Math & 85 \tabularnewline
Statistics & 81 \tabularnewline
Computing & 35 \tabularnewline
Business & 29 \tabularnewline
\bottomrule
\end{longtable}

In 2008, this wasn't a math blog---it was a consultant's notebook, a professional sharing whatever seemed useful that day.

\bookfigure{ch01_activity_heatmap.png}{Activity heatmap, 2008-2025}

\bigskip

\section*{The First Posts}

What a working consultant found interesting in 2008:

\textbf{January 9:} "Moore's law and software bloat": Software expands to fill available hardware.

\textbf{January 11:} "Million dollar cutoff for software technique": A practical threshold for when formal methods pay off.

\textbf{January 14:} "Tips for learning regular expressions": The tools of the trade.

\textbf{January 21:} "Selection bias and bombers": A curiosity, a calculation, a small insight.

These were not the posts of someone building a platform, but of someone thinking out loud---a helpful but tentative voice making no claim to authority.

Worth quoting the very first post in full:

\begin{blogquote}
I ran across an article recently comparing the performance of a 1986 Mac and a 2007 PC. Of course the new machine would totally blow away the old one on a number crunching benchmark, but when it comes to the most mundane benchmarks --- time to boot, launch Microsoft Word, open a file, do a search and replace, etc. --- the old Mac pulls ahead slightly. Software bloat has increased at roughly the same rate as Moore's law, making a new machine with new software no better than an old machine with old software in some respects.
\end{blogquote}

\begin{blogquote}
The comparisons in the article resonate with my experience. I expect administrative tasks to be quick and number crunching to be slow, and so I'm continually surprised how long routine tasks take and how quickly numerical software runs.
\end{blogquote}

\begin{blogquote}
Update (December 2014): Processors have not gotten much faster in the time since this post was first published, though we have more of them. However, applications do open faster. Perhaps software vendors have paid attention to bloat precisely because the free lunch of faster processors is over.
\end{blogquote}

One hundred eighty-two words that started a seventeen-year practice.

\bigskip

\section*{The Consultant's Breadth}

In 2008, Cook was a research statistician at MD Anderson Cancer Center, consulting on clinical trial design. But he'd also been a software developer, a project manager, and an independent consultant---someone who worked on many things, with breadth reflecting that history.

Clinical trials appear in 55 posts over the blog's lifetime, concentrated in the early years. By 2018, they had faded to a single post. The topic was professional obligation, not personal fascination.

PowerShell appears in 47 posts: the tool of a consultant working in Windows environments. Cook wrote about PowerShell from 2008 to 2012, with the heaviest concentration in the first two years. Then PowerShell fades, replaced by Python.

Perl gets 20 posts scattered across 2008-2021, with the last appearing in March 2021.

These are fossils, preserving what a working consultant needed to know in the late 2000s. Their fading reveals temporary interest versus permanent fascination.

\bigskip

\section*{The Early Rhythm}

Even in 2008, patterns emerge.

\textbf{Tuesday} leads with 86 posts. \textbf{Wednesday} and \textbf{Thursday} follow with 65 each. Weekends trail: Saturday has 48, Sunday just 35.

This is a working professional's rhythm---attention when work allows, weekday mornings, before the demands take over. Already visible in year one: the 7 AM peak that persists through all seventeen years.

\textbf{Morning} dominates. In 2008, the hour from 7:00-7:59 AM accounts for a little over 10\% of posts; across the full corpus it reaches 13\%. Add the hours from 6-9 AM and you have about 30\% of 2008's output and 43\% of all posts over seventeen years.

Cook was writing before work. Or perhaps writing \textit{was} work: a way to warm up the mind, to think through problems, to share what seemed interesting before the meetings started.

\bigskip

\section*{The Short-Post Era}

The median post length in 2008 was 163 words---half of all posts were shorter than that.

Some were mere announcements: "New prime number record" (9 words). Others were pointers: "Technology hype curve" (12 words). The blog's \textbf{median} post length overall is 312 words, but in 2008 it was just 163. The blog served a curatorial function in 2008, directing attention more than providing explanation.

The deep dives would come later. First came the practice establishing itself---the habit of showing up, of publishing something, anything, every day.

Quality would follow quantity, but the habit of quantity had to be established first.

\bigskip

\section*{What 2008 Reveals}

What beginning looks like:

\textbf{Scattered.} Topics ranging from software bloat to Chinese Unicode to clinical trials to PowerShell. No focus because focus hasn't been found yet.

\textbf{Short.} Two hundred six words on average: enough to make a point, not enough to develop an argument. Notes, not essays.

\textbf{Frequent.} More than one post per day. Beginning's energy, before sustainability matters.

\textbf{Professional.} Topics reflecting a consultant's working life. Serving the job.

Many new blogs start this way. Few sustain it. Typical pattern: a burst of enthusiasm, then gradual abandonment. The internet littered with blogs that have three posts from 2009 and then nothing.

Cook's first year was not a burst that would fade---it was a beginning.

\bigskip

\section*{The Years That Followed}

2008 established the habit. Following years tested it.

\begin{longtable}{p{3.00cm}p{3.00cm}p{3.00cm}}
\toprule
\textbf{Year} & \textbf{Posts} & \textbf{Avg Words} \tabularnewline
\midrule
2008 & 410 & 206 \tabularnewline
2009 & 366 & 256 \tabularnewline
2010 & 284 & 268 \tabularnewline
2011 & 317 & 222 \tabularnewline
2012 & 357 & 216 \tabularnewline
\bottomrule
\end{longtable}

Post count declined from 410 to 284-366, then stabilized. Average length crept up from 206 to 256-268. The trade-off between quantity and depth was beginning: more words per post, fewer posts per year.

But the habit held. Three hundred or more posts every year. Practice taking root.

\bookfigure{ch01_posts_wordcount.png}{Posts and length by year}

\bigskip

\section*{The Notebook Phase}

Years 2008-2012 form the first era. Call it the Notebook Phase.

Total posts: 1,734
Average length: 232 words
Dominant categories: Software development, Math, Computing, Business
Characteristic posts: Tips, pointers, short observations

A consultant's scratchpad. Useful observations from working life. Tools and techniques. Quick notes on mathematical curiosities. Nothing claiming to be comprehensive or authoritative.

The voice was helpful but not yet distinctive, the topics reflected professional needs, and the posts served the moment rather than building toward permanence.

Remarkable not for quality but consistency. Five years of showing up. Five years of writing something, even when there was nothing profound to say. Five years building the habit that would sustain what came next.

Start before you're ready. The 2008 blog wasn't waiting for quality; it was generating quantity until quality emerged. Four hundred ten posts of variable merit, but four hundred ten chances to find what worked.


